\chapter{The Hawking Effect}
\label{ch:curved-hawking-effect}

\ConventionsLorentzian

The Hawking effect is a statement about QFT on a fixed spacetime containing a
black-hole horizon.  It relates regularity of a state at the future horizon
to a thermal flux at future null infinity.  The temperature is determined by
the surface gravity of the horizon.

\section{Stationary Horizon and Surface Gravity}

Let \((M,g)\) be a stationary black-hole spacetime with Killing vector
\(\chi^\mu\) generating the horizon.  On the horizon,
\[
  \chi^\nu\nabla_\nu\chi^\mu=\kappa\chi^\mu ,
\]
where \(\kappa\) is the surface gravity.  Near a nonextremal horizon, the
metric in the two-dimensional normal directions has Rindler form
\[
  ds^2\simeq -\kappa^2\rho^2dt^2+d\rho^2+\cdots .
\]
Thus the local horizon problem is the Rindler problem of
Chapter~\ref{ch:curved-unruh-effect-modular-theory}, with acceleration scale
replaced by \(\kappa\).

\section{Euclidean Regularity}

Analytically continue \(t=-\ii\tau\).  The normal metric becomes
\[
  ds_E^2\simeq \kappa^2\rho^2d\tau^2+d\rho^2+\cdots .
\]
This is smooth at \(\rho=0\) only if the angular variable
\(\kappa\tau\) has period \(2\pi\).  Therefore the Euclidean state regular at
the horizon is periodic in imaginary Killing time with inverse temperature
\[
  \beta_H=\frac{2\pi}{\kappa},
  \qquad
  T_H=\frac{\kappa}{2\pi}.
\]
This derivation establishes the KMS periodicity of correlation functions that
admit the Euclidean regular continuation; it does not by itself prove
existence of the interacting Lorentzian state.

\section{Bogoliubov Coefficients From Mode Tracing}

The Euclidean argument identifies the temperature of a regular stationary
state.  The collapse calculation gives a Lorentzian particle-production
mechanism.  We now carry out the coefficient calculation for a massless
scalar field in the geometric-optics approximation.  The angular variables
and the potential barrier are suppressed in this section; they are restored
as greybody factors below.

Let \(v\) be advanced time on past null infinity \(\mathcal I^-\), and let
\(u\) be retarded time on future null infinity \(\mathcal I^+\).  In a
collapsing geometry forming a nonextremal future horizon, outgoing null rays
arriving at late \(u\) originate from \(v\) exponentially close to the last
escaping ray \(v_0\).  The relation has the universal form
\[
  v_0-v=C e^{-\kappa u}+O(e^{-2\kappa u}),
  \qquad C>0.
\]
Equivalently,
\[
  u(v)=-\frac{1}{\kappa}\log\frac{v_0-v}{C}+O(v_0-v),
  \qquad v<v_0 .
\]
The constant \(C\) depends on the collapse geometry and on the normalization
of \(u\), but it will only contribute a phase to the occupation number.

Normalize right-moving modes on a null line by
\[
  f_{\omega'}^{\rm in}(v)
  =
  \frac{1}{\sqrt{4\pi\omega'}}e^{-\ii\omega' v},
  \qquad
  p_{\omega}^{\rm out}(u)
  =
  \frac{1}{\sqrt{4\pi\omega}}e^{-\ii\omega u},
  \qquad
  \omega,\omega'>0 .
\]
Tracing \(p_\omega^{\rm out}\) backwards to \(\mathcal I^-\) gives, near the
late-time singular part responsible for particle creation,
\[
  p_\omega(v)
  =
  \frac{1}{\sqrt{4\pi\omega}}
  \left(\frac{v_0-v}{C}\right)^{\ii\omega/\kappa}\theta(v_0-v)
  +\hbox{terms smooth across }v=v_0 .
\]
The smooth terms have rapidly decaying Fourier transform at large
\(\omega'\) and do not determine the thermal ratio.  The singular term is
expanded as
\[
  p_\omega(v)
  =
  \int_0^\infty d\omega'\,
  \left(
  \alpha_{\omega\omega'}f_{\omega'}^{\rm in}(v)
  +
  \beta_{\omega\omega'}\overline{f_{\omega'}^{\rm in}(v)}
  \right).
\]
Using the Klein-Gordon product on a null line,
\[
  (f,g)_{\mathcal I^-}
  =
  -\ii\int_{-\infty}^{\infty}
  f(v)\overleftrightarrow{\partial_v}\overline{g(v)}\,dv,
\]
one obtains
\[
  \alpha_{\omega\omega'}=(f_{\omega'}^{\rm in},p_\omega)_{\mathcal I^-},
  \qquad
  \beta_{\omega\omega'}=-(\overline{f_{\omega'}^{\rm in}},p_\omega)_{\mathcal I^-}.
\]
Up to phases depending on \(C\) and \(v_0\), the required Fourier integrals
are
\[
  \int_0^\infty x^{\ii a}e^{-\ii\omega' x}\,dx
  =
  e^{-\ii\pi/2}e^{\pi a/2}\Gamma(1+\ii a)(\omega')^{-1-\ii a},
\]
\[
  \int_0^\infty x^{\ii a}e^{+\ii\omega' x}\,dx
  =
  e^{+\ii\pi/2}e^{-\pi a/2}\Gamma(1+\ii a)(\omega')^{-1-\ii a},
  \qquad a=\frac{\omega}{\kappa}.
\]
Equivalently, since \(\Gamma(1+\ii a)=\ii a\,\Gamma(\ii a)\), the singular
Bogoliubov coefficients can be written in the form
\[
  \alpha_{\omega\omega'}
  =
  A_{\omega\omega'}\,e^{+\pi\omega/(2\kappa)}
  \Gamma(\ii\omega/\kappa)(\omega')^{-\ii\omega/\kappa},
\]
\[
  \beta_{\omega\omega'}
  =
  -A'_{\omega\omega'}\,e^{-\pi\omega/(2\kappa)}
  \Gamma(\ii\omega/\kappa)(\omega')^{-\ii\omega/\kappa},
\]
where \(A_{\omega\omega'}\) and \(A'_{\omega\omega'}\) differ only by phases
and have common modulus
\[
  |A_{\omega\omega'}|
  =
  \frac{1}{2\pi\kappa}\sqrt{\frac{\omega}{\omega'}} .
\]
Therefore
\[
  \frac{|\alpha_{\omega\omega'}|^2}{|\beta_{\omega\omega'}|^2}
  =
  e^{2\pi\omega/\kappa}.
\]
Using
\[
  |\Gamma(\ii a)|^2=\frac{\pi}{a\sinh(\pi a)}
\]
gives the explicit negative-frequency coefficient
\[
  |\beta_{\omega\omega'}|^2
  =
  \frac{1}{2\pi\kappa\,\omega'}
  \frac{1}{e^{2\pi\omega/\kappa}-1}.
\]
The factor \(1/\omega'\) is not a typographical nuisance; it records that the
idealized late-time stationary approximation has infinite retarded-time
duration.  A mathematically meaningful particle number is obtained by using
wave packets localized in retarded time.  For wave packets centered at
frequency \(\omega_j\) and late retarded-time bin \(n\), the expectation
value tends to
\[
  \langle N_{j n}\rangle
  =
  \frac{1}{e^{2\pi\omega_j/\kappa}-1}
  +\hbox{terms vanishing at late }n ,
\]
before greybody scattering is included.  Thus
\[
  n_\omega=\frac{1}{e^{\beta_H\omega}-1},
  \qquad
  \beta_H=\frac{2\pi}{\kappa}.
\]

The same tracing relation shows the trans-Planckian sensitivity of the
geometric-optics derivation.  A wave packet observed at retarded time \(u\)
with frequency \(\omega\) has precursor frequency on \(\mathcal I^-\) of
order
\[
  \omega_{\rm in}\sim
  \omega\,\frac{du}{dv}
  =
  \frac{\omega}{\kappa(v_0-v)}
  \sim
  \frac{\omega}{\kappa C}e^{\kappa u}.
\]
At sufficiently late \(u\) this exceeds any fixed microscopic scale.  The
fixed-background QFT derivation is therefore a conditional theorem inside
the assumed continuum theory; robustness under modified short-distance
physics is a separate problem, often modeled in two dimensions by moving
mirror trajectories \(v=p(u)=v_0-Ce^{-\kappa u}\), which reproduce the same
Bogoliubov ratio and stress-tensor flux for a conformal scalar.

\section{States: Boulware, Hartle--Hawking, and Unruh}

The Boulware state is empty with respect to static time at infinity but is
singular at the horizon.  The Hartle--Hawking state is regular on both future
and past horizons and is thermal on both exterior regions.  The Unruh state
is regular on the future horizon and empty on past null infinity; it describes
an evaporating black hole in the fixed-background approximation and produces
outgoing thermal flux at future null infinity.

\section{Flux and Greybody Factors}

In an asymptotically flat black-hole spacetime, propagation from the near
horizon region to infinity includes scattering by the curvature potential.
The number flux for angular momentum channel \(\ell\) has the form
\[
  \frac{dN_{\ell}}{dt\,d\omega}
  =
  \frac{\Gamma_\ell(\omega)}
  {2\pi\bigl(e^{\beta_H\omega}-1\bigr)},
\]
where \(\Gamma_\ell(\omega)\) is the greybody factor determined by the
classical wave equation in the exterior geometry.  The thermal factor is
horizon kinematics; the greybody factor is propagation through the exterior
potential.

\section{QFT Boundary}

For interacting QFT, a theorem-level Hawking effect requires existence of a
Hadamard state with the required horizon regularity, control of stress-tensor
renormalization, and propagation estimates from the horizon to infinity.  In
semiclassical gravity one must also solve the backreaction equation
\[
  G_{\mu\nu}=8\pi G\,\langle T_{\mu\nu}\rangle_{\rm ren},
\]
which is beyond fixed-background QFT.

\begin{openproblem}[Interacting Hawking theorem]
\label{op:interacting-hawking-theorem}
Prove a Hawking radiation theorem for an interacting local QFT on a class of
black-hole backgrounds, including Hadamard state construction, KMS or
near-horizon analyticity, renormalized stress tensor, greybody propagation,
and the precise limits in which the emitted flux is thermal.
\end{openproblem}
